\documentclass[12pt,a4paper]{article}
\usepackage[utf8]{inputenc}
\usepackage[T1]{fontenc}
\usepackage[brazil]{babel}
\usepackage{geometry}
\usepackage{setspace}
\usepackage{booktabs}
\usepackage{siunitx}
\usepackage{hyperref}
\geometry{margin=2.5cm}
\onehalfspacing

\title{Relat\'orio Acad\^emico\\Simula\c{c}\~ao de Risco para Custo de Constru\c{c}\~ao de Gasoduto}
\author{Grupo 6\\Bruna Irene da Silva (119062418)\\Davi dos Santos Mattos (119133049)\\Gabriel Mayrink Verdun (125078427)}
\date{\today}

\begin{document}

\maketitle

\section{Introdu\c{c}\~ao}
Este relat\'orio apresenta a an\'alise de risco do custo de um projeto de constru\c{c}\~ao de gasoduto, com base em simula\c{c}\~ao de Monte Carlo implementada em R. O objetivo \`e quantificar a incerteza do custo total, estimar um valor-base para or\c{c}amento, propor uma margem de conting\^encia e avaliar a viabilidade de uma proposta comercial de \$45 milh\~oes.

\section{Metodologia}
Foram realizadas 3000 simula\c{c}\~oes, modelando incertezas de rota e de componentes de custo.

\subsection{Estrutura do modelo}
A rota do gasoduto foi modelada como vari\'avel aleat\'oria bin\'aria:
\begin{itemize}
  \item Rota de 260 km;
  \item Rota alternativa de 290 km.
\end{itemize}
A probabilidade da rota alternativa foi sorteada de forma uniforme no intervalo [0,35; 0,45], incorporando incerteza de cen\'ario no tra\c{c}ado final.

\subsection{Vari\'aveis de custo}
Os principais insumos econ\^omicos foram modelados por distribui\c{c}\~oes triangulares:
\begin{itemize}
  \item custo de tubo,
  \item tempo de escava\c{c}\~ao,
  \item custo de m\~ao de obra,
  \item transporte,
  \item soldagem,
  \item filtro,
  \item acabamento por quil\^ometro.
\end{itemize}
O custo total da obra foi obtido como soma de tr\^es blocos: material, m\~ao de obra e servi\c{c}os.

\section{Resultados}
A simula\c{c}\~ao gerou a distribui\c{c}\~ao emp\'irica do custo total (em milh\~oes de USD), analisada por histogramas, densidade e fun\c{c}\~ao de distribui\c{c}\~ao acumulada.

\subsection{M\'etricas de refer\^encia}
Os resultados centrais reportados no notebook foram:

\begin{table}[h!]
\centering
\begin{tabular}{lr}
\toprule
M\'etrica & Valor (USD milh\~oes) \\
\midrule
Custo base (P50) & 44,52 \\
Pre\c{c}o proposto de refer\^encia (P80) & 47,78 \\
Conting\^encia absoluta & 3,26 \\
Conting\^encia relativa & 7,32\% \\
\bottomrule
\end{tabular}
\caption{Resumo dos indicadores de or\c{c}amento e conting\^encia.}
\end{table}

\subsection{An\'alise da proposta de USD 45 milh\~oes}
A proposta de USD 45 milh\~oes foi comparada com a distribui\c{c}\~ao simulada e com os percentis de refer\^encia. No notebook, a conclus\~ao final foi de \textbf{n\~ao aceita\c{c}\~ao} da proposta, por margem insuficiente frente ao risco de estouro de custo.

\section{Discuss\~ao}
A diferen\c{c}a entre o P50 e o P80 indica assimetria e dispers\~ao relevantes na distribui\c{c}\~ao de custos. Em termos gerenciais:
\begin{itemize}
  \item utilizar o P50 como base representa postura neutra de risco;
  \item utilizar o P80 como pre\c{c}o-alvo incorpora robustez financeira;
  \item uma proposta pr\'oxima ao P50 tende a elevar a probabilidade de perda econ\^omica para o contratante.
\end{itemize}
Assim, o valor de USD 45 milh\~oes n\~ao oferece folga adequada para absorver variabilidades operacionais e de produtividade modeladas na simula\c{c}\~ao.

\section{Conclus\~ao}
A modelagem estoc\'astica adotada foi adequada para suportar decis\~ao or\c{c}ament\'aria em condi\c{c}\~oes de incerteza. Com base nos resultados:
\begin{itemize}
  \item custo-base recomendado: aproximadamente USD 44,52 milh\~oes (P50);
  \item valor com maior seguran\c{c}a: aproximadamente USD 47,78 milh\~oes (P80);
  \item conting\^encia recomendada: USD 3,26 milh\~oes (7,32\%);
  \item proposta de USD 45 milh\~oes: classificada como arriscada e n\~ao recomendada.
\end{itemize}

\section{Limita\c{c}\~oes e trabalhos futuros}
Como extens\~ao, recomenda-se: (i) calibra\c{c}\~ao dos par\^ametros com dados hist\'oricos de obras similares; (ii) an\'alise de sensibilidade para identificar vari\'aveis de maior impacto; e (iii) incorpora\c{c}\~ao de cen\'arios macroecon\^omicos (infla\c{c}\~ao, c\^ambio e juros) no modelo de custos.

\end{document}
